\documentclass[12pt,a4paper]{article}
\usepackage[utf8]{inputenc}
\usepackage[english]{babel}
\usepackage{graphicx}
\usepackage{amsmath}
\usepackage{amssymb}
\usepackage{hyperref}
\usepackage{booktabs}
\usepackage{geometry}
\usepackage{natbib}
\usepackage{float}
\usepackage{url}

\geometry{a4paper, margin=1in}

\title{Stock Market Prediction Using Historical Data}
\author{Your Name}
\date{\today}

\begin{document}

\maketitle

\begin{abstract}
This report presents a machine learning approach to predict stock market movements based on historical price data. Multiple models including traditional time series methods and modern machine learning techniques are implemented and evaluated. The results show that [to be completed with findings]. The report discusses the methodology, implementation, results, and potential applications of the developed prediction system.
\end{abstract}

\tableofcontents
\newpage

\section{Introduction}
\subsection{Background and Motivation}
Stock market prediction is a challenging problem that has attracted attention from researchers in finance, statistics, and machine learning. The ability to predict market movements, even with modest accuracy, can potentially lead to profitable trading strategies and better risk management.

\subsection{Problem Statement}
This project aims to develop a machine learning model that can predict stock prices or movement directions using historical market data. The specific objectives include:
\begin{itemize}
    \item Evaluate the effectiveness of different machine learning algorithms for stock price prediction
    \item Identify useful technical indicators and features for prediction
    \item Compare performance across different time horizons and market conditions
    \item Assess the practical applicability of such predictions
\end{itemize}

\subsection{Approach Overview}
The project follows a standard machine learning workflow including data preprocessing, feature engineering, model development, and evaluation. Multiple models are trained and compared to identify the most effective approach.

\section{Literature Review}
[Insert summary of literature review here based on literature\_review.md]

\section{Data and Methodology}
\subsection{Dataset Description}
The dataset used in this study consists of historical stock market data, including opening and closing prices, high and low prices, and trading volumes. The data covers [time period] for [market/stocks].

\subsection{Data Preprocessing}
Before training the models, several preprocessing steps were applied to the raw data:
\begin{itemize}
    \item Handling missing values
    \item Normalization and scaling
    \item Feature engineering
    \item Train/test splitting
\end{itemize}

\subsection{Feature Engineering}
Several technical indicators were calculated and used as features:
\begin{itemize}
    \item Moving averages (5-day, 20-day)
    \item Relative Strength Index (RSI)
    \item Moving Average Convergence Divergence (MACD)
    \item Volatility measures
    \item Price momentum indicators
\end{itemize}

\subsection{Models}
The following models were implemented and evaluated:
\begin{itemize}
    \item Linear regression with regularization
    \item Random Forest
    \item Gradient Boosting
    \item Support Vector Regression
\end{itemize}

\section{Experimental Setup}
\subsection{Training and Testing Methodology}
The models were trained on historical data up to [date] and tested on subsequent data. A walk-forward validation approach was used to simulate real-world trading conditions.

\subsection{Evaluation Metrics}
The models were evaluated using the following metrics:
\begin{itemize}
    \item Mean Squared Error (MSE)
    \item Root Mean Squared Error (RMSE)
    \item Mean Absolute Error (MAE)
    \item Direction Accuracy
    \item R-squared
\end{itemize}

\section{Results and Discussion}
\subsection{Model Performance Comparison}
[Insert table or figures comparing model performance]

\subsection{Feature Importance Analysis}
[Insert analysis of most important features]

\subsection{Error Analysis}
[Insert analysis of where models performed well or poorly]

\subsection{Discussion}
The experimental results show [discuss key findings]. These findings align with/contradict previous research that [discuss relationship to literature]. The implications for trading and investment strategies are [discuss practical implications].

\section{Conclusion and Future Work}
\subsection{Summary of Findings}
[Summarize key findings and contributions]

\subsection{Limitations}
The current approach has several limitations:
\begin{itemize}
    \item Limited to technical indicators without fundamental or sentiment data
    \item May not account for market regime changes
    \item Potential issues with market efficiency and predictability
\end{itemize}

\subsection{Future Work}
Potential directions for future work include:
\begin{itemize}
    \item Incorporating news sentiment and alternative data
    \item Developing ensemble methods combining multiple models
    \item Implementing deep learning approaches like LSTM networks
    \item Extending to portfolio optimization and risk management
\end{itemize}

\section{References}
% Use the natbib package for references
\bibliographystyle{plainnat}
\bibliography{references}

\appendix
\section{Additional Results}
[Insert any additional tables, figures, or analysis]

\section{Code and Implementation Details}
Details of the implementation, including key algorithms and code snippets, are available in the project repository: [link to repository].

\end{document} 